\section{Discussion}
\label{sec:discussion}

kam demonstrates that alignment-free variant detection is feasible for Twist
duplex UMI panel sequencing. The pipeline preserves molecule-level information
from raw FASTQ to variant call without any alignment step, and detects
40--48\% of truth variants at 2\% VAF with precision above 0.60. The pipeline
runs end-to-end in under 25 seconds per sample on a single core at 200K read
pairs.

The primary result is proof of concept: k-mer-based de Bruijn graph traversal
can recover somatic variants from duplex UMI panel data without reference
alignment. The sensitivity numbers reported here are not representative of the
pipeline's upper bound; they reflect two implementation bugs (duplex orientation
comparison) and one hardware constraint (200K read limit for memory safety),
both of which are tractable to fix.

\paragraph{Comparison to alignment-based approaches.}
Alignment-based pipelines for ctDNA variant calling typically report sensitivity
of 80--95\% at 1--2\% VAF with comparable panel sizes, using millions of
read pairs~\autocite{Twist22cfdna}. The gap between those numbers and the
40--48\% reported here is entirely explained by depth (200K vs.\ $>$1M reads)
and by the loss of duplex signal. At matching depth with duplex orientation
corrected, we expect sensitivity to increase substantially, though a full head-to-head comparison is left to future work.

\paragraph{Information preservation.}
The key architectural advantage of kam over the km/Jellyfish pipeline is that
every k-mer in the graph retains a pointer to the molecule that contributed it.
This means strand bias, duplex support, and per-molecule error probability are
available at the scoring stage without a separate alignment pass. Current calling
uses only molecule counts and strand counts; the full posterior over molecule
quality (incorporating per-cycle error probability from base quality scores)
is implemented but not yet tuned.

\paragraph{Reproducibility.}
All RNG in the pipeline uses seeded deterministic generators, ensuring that
runs on the same input produce identical output. This is a requirement for
Nextflow cache compatibility and for scientific reproducibility. The pipeline
does not use the system clock, process ID, or any other non-deterministic
seed source.
